
\lesson{30/10}
[...time evolution operator and interacting field, sec 8.3.1 notes]
In the interacting picture \textbf{the states evolve in time according to the
interaction Hamiltonian, while operators evolve according to the free Hamiltonian
.} Time dependency has a \textit{reference time} $t_0$, in which the regular QM
Heisenberg's picture coincide with the interaction picture
(i.e. before the perturbation starts).

The field in the interaction picture is given by
$$\phi_I(t,\mathbf{x}):= e^{iH_0(t-t_0)}\phi(t_0,\mathbf{x})e^{-iH_0(t-t_0)}$$
and all other operators are similarly defined, like the interaction Hamiltonian
$$H_I:=e^{iH_0(t-t_0)}H_\text{int}e^{-iH_0(t-t_0)}$$
where $H_\text{int}$ is built with the fields $\phi(t_0,\mathbf{x})$.

Inverting the relation between $\phi,\phi_I$ yields
$$\phi(t,\mathbf{x}):= e^{-iH_0(t-t_0)}\phi_I(t_0,\mathbf{x})e^{iH_0(t-t_0)}.$$
\textbf{Note} The interacting field $\phi_I$ is a \textit{free field}, because
it is the time evolution of a free field with a free Hamiltonian: it can be
interpreted as the evolution of the initial condition $\phi(t_0,\mathbf{x})$ with
the interacting Hamiltonian switched off.

The full field $\phi(t,\mathbf{x})$ is defined as
$$\phi(t,\mathbf{x})=e^{iH(t-t_0)}\phi(t_0,\mathbf{x})e^{-iH(t-t_0)}\footnote{
non sono sicuro, se H(t) posso usare la forma exp????}$$
inserting the expression of $\phi(t_0)$ in terms of $\phi_I(t)$ yields
$$\phi(t,\mathbf{x})=\underbrace{e^{iH(t-t_0)}e^{-iH_0(t-t_0)}}_{
:=U_I^\dagger(t,t_0)} \phi_I(t,\mathbf{x}) \underbrace{e^{iH_0(t-t_0)}
e^{-iH(t-t_0)}}_{:=U_I(t,t_0)}
$$
where $U_I$ is the \textbf{time evolution operator in the interacting picture}.

The derivative of $U_I$ is
$$i\pdv[]{}{t}U_I = e^{iH_0(t-t_0)}(H-H_0)e^{-iH(t-t_0)}$$
recalling $H-H_0=H_\text{int}$ and inserting
$\mathbb{I}=e^{=iH_0(t-t_0)}e^{iH_0(t-t_0)}$ yields
$$i\pdv[]{}{t}U_I(t,t_0) = \underbrace{e^{iH_0(t-t_0)}H_\text{int}
e^{-iH_0(t-t_0)}}_{H_I(t)}\underbrace{e^{iH_0(t-t_0)}e^{-iH(t-t_0)}}_
{U_I(t,t_0)}$$
So finally
$$i\pdv[]{}{t}U_I(t,t_0) = H_I(t)U_I(t,t_0)$$
This differential equation is not trivial to solve as in QM, since $H$ is a
time dependent operator and causality may become an issue, since
$$[H_I(t_1),H_I(t_2)]=0$$
\textbf{does not hold} $\mathbf{\forall t_1,t_2}$, so the time ordering operator
is here needed.

The QM form of the solution is still correct:
$$U_I(t,t_0) = T\exp\left[-i\int_{t_0}^t \dd t' H_I(t')\right]=
\mathbb{I} + (-i)\int_{t_0}^t \dd t' H_I(t') + \frac{(-i)^2}{2!}\int_{t_0}^t
\dd t_1\int_{t_0}^{t}\dd t_2 T\left[H_I(t_1)H_I(t_2)\right]+...
$$
but a problem arises with $\pdv[]{}{t}U_I(t,t_0)$, since the double (and higher
order) products' derivatives have problems with commutators:
$$\pdv[]{}{t}\left(\int_{t_0}^t \dd t_1 H_I(t_1) \int_{t_0}^t \dd t_2 H_I(t_2)
\right)=H_I(t)\int_{t_0}^t \dd t'' H_I(t'') + \int_{t_0}^t \dd t'H_I(t')$$
to reorder the second term the commutator between $H_I(t)$ and $H_I(t')$ must be
null. To guarantee this, the \textit{time ordering operator} is put at the
start of the series expansion\\
\textbf{!!!RIVEDI RIFAI ETC!!! PROVLEMA CON 2! ed estremi e T}\\
where $t_0<t_2<t_1<t$ is guaranteed by $T$.\\
\COMMENT{\textbf{Note} The factor $\flatfrac{1}{2!}$ comes from the fact that the double
integral's original form is
$$(-i)^2\int_{t_0}^t \dd t_1\int_{t_0}^{t_1} \dd t_2 H_I(t_1)H_I(t_2)$$
and so the price to pay for having the same extrema in $t_1$ and $t_2$ is a factor
$\flatfrac{1}{2}$, which can be generalized for the $n-$th term to
$\flatfrac{1}{n!}$.}

Two useful properties of the time evolution operator are
$$U_I(t_1,t_2)U_I(t_2,t_3)=U_I(t_1,t_3),\,\forall t_1\geq t_2\geq t_3$$
$$U_I(t_1,t_3)U_I^\dagger(t_2,t_3) = U_I(t_1,t_2),\,\forall t_1\geq t_2\geq t_3$$
\subsection{Ground and vacuum}
%You shall seal the blinding light that plagues their dreams
The vacuum state is so that
$$H_0\ket{0}=0$$
while the ground state of an interacting theory
$$H\ket{\Omega}=E_0\ket{\Omega}$$
with, in general, $E_0\neq 0$.

Consider the complete set of eigenstates of the full Hamiltonian
$$\ket{n}\to\sum_{n}^{}\ket{n}\bra{n}=\mathbb{I}$$
and evolve the vacuum state up to a time $T>0$
$$e^{-iHT}\ket{0} = e^{-iHT}\sum_{n}^{}\ket{n}\bra{n}\ket{0}=
\sum_{n}^{}e^{-iE_nT}\ket{n}\bra{n}\ket{0}
$$
since $\ket{\Omega}=\ket{n=0}$ and the interaction is small,
there is an overlap of $\ket{0}$ and $\ket{\Omega}$:
$$e^{-iHT}\ket{0} = e^{-iE_0T}\ket{\Omega}\bra{\Omega}\ket{0} + \sum_{n\neq0}^{}
e^{-iE_nT}\ket{n}\bra{n}\ket{0}.$$
In the limit $T\to\infty$, the first term is dominant, since $E_0\leq E_n$. The
formal way to do this is to "tilt" the limit of $T\to\infty$ to
$$T\to(\infty(1-i\varepsilon))$$
And so taking the limit, neglecting $n>0$ terms and inverting the last relation
$$\ket{\Omega} = \lim\limits_{T\to\infty(1-i\varepsilon)}\left(e^{-iE_0T}\bra{\Omega}
\ket{0}\right)^{-1}e^{-iHT}\ket{0}$$
Since $T$ goes to $\infty$, adding a constant to it does not change anything:
$$T\to T+t_0 = T-(-t_0)$$
So
$$\ket{\Omega} = \lim\limits_{T\to\infty(1-i\varepsilon)}\left(e^{-iE_0(T-(-t_0)}
\bra{\Omega}\ket{0}\right)^{-1}e^{-iH(T-(-t_0))}\ket{0}$$
substituting now
$$\ket{0}=e^{iH_0(t_0-(-T))}\ket{0}$$
yields
$$\ket{\Omega}=\lim\limits_{T\to\infty(1-i\varepsilon)}\left(e^{-iE_0(T-(-t_0)}
\bra{\Omega}\ket{0}\right)^{-1}e^{-iH(T-(-t_0))}e^{iH_0(t_0-(-T))}\ket{0}$$
so finally
$$\boxed{
\ket{\Omega}=\lim\limits_{T\to\infty(1-i\varepsilon)}\left(e^{-iE_0(t_0-(-T)}
\bra{\Omega}\ket{0}\right)^{-1}U_I(t_0,-T)\ket{0}}$$
which means that the ground state $\ket{\Omega}$ can be found by evolving the
vacuum state.
$$\boxed{
\bra{\Omega} = \lim\limits_{T\to\infty(1-i\varepsilon)}\bra{0}U_I(T,t_0)
\left(e^{-iE_0(T-t_0)}
\bra{\Omega}\ket{0}\right)^{-1}
}$$\\

\subsection{Gell-Mann-Low Formula}
So now the full propagator formula can be derived from its definition
[FORMATTAZIONE]

\begin{align*}
\bra{\Omega}\phi(x)\phi(y)\ket{\Omega}=\lim\limits_{T\to\infty(1-i\varepsilon)}
\left(e^{-iE_0(T-t_0)}\bra{\Omega}\ket{0}\right)^{-1}&
\bra{0}U_I(T,t_0) U_I^\dagger(x^0,t_0)\phi_I(x)U_I(x^0,t_0)U^\dagger_I(y^0,t_0)
\phi_I(y)\\&
U_I(y^0,t_0)U_I(t_0,-T)\ket{0}\left(e^{-iE_0(t_0-(-T)}
\bra{\Omega}\ket{0}\right)^{-1}
\end{align*}
by the properties of the time evolution operator, some products collapse and
the constants can be simplified by noting
$$1=\bra{\Omega}\ket{\Omega} = \left(\left|\bra{0}\ket{\Omega}\right|^2e^{-iE_02T
}\right)^{-1}\bra{0}U_I(T,t_0)U_I(t_0,-T)\ket{0}$$
from which
$$\left(\left|\bra{\Omega}\ket{0}\right|^2e^{-iE_02T}\right)^{-1}=\frac{1}{
\bra{0}U_I(T,-T)\ket{0}}.$$
The final expression, also called the \textit{reduction formula}, is
$$\boxed{\bra{\Omega}T\phi(x)\phi(y)\ket{\Omega} = 
\lim\limits_{T\to\infty(1-i\varepsilon)}\frac{\bra{0}T\left\{\phi_I(x)\phi_I(y)
\exp\left[-i\int_{-T}^T \dd t H_I(t)\right]\right\}\ket{0}}
{\bra{0}T\left\{\exp\left[-i\int_{-T}^T \dd t H_I(t)\right]\right\}\ket{0}}}
$$
This formula has an elegant expression in terms of the $S$-matrix
$$S:=\lim\limits_{T\to\infty(1-i\varepsilon)} T\left\{\exp\left[-i\int_{-T}^T \dd t
H_I(t)\right]\right\}$$
which is
$$\boxed{\bra{\Omega}T\phi(x)\phi(y)\ket{\Omega} = \frac{\bra{0}T\{\phi_I(x)
\phi_I(y)S\}\ket{0}}{\bra{0}S\ket{0}}}$$
\subsection{Wick's Theorem}
This theorem is very important, because it relates \textit{time ordering}
with \textit{normal ordering} (which is easier to compute on the vacuum state).

The first step is to define two separate parts of a scalar field
$$\phi(x)=\phi^+(x)+\phi^-(x)$$
where
$$\phi^-(x)\propto \int\dd^3k a_k,\;\phi^+(x)\propto\int \dd^3ka^\dagger_k.$$
So the product
$$\phi(x)\phi(y) = \phi^+(x)\phi^+(y) + \underbrace{\phi^-(x)\phi^+(y)}_\text{not
ordered}+\phi^+(x)\phi^-(y)+\phi^-(x)\phi^-(y)$$
is almost normal ordered
$$N(\phi(x)\phi(y))=\phi^+(x)\phi^+(y) + \phi^+(y)\phi^-(x)
+\phi^+(x)\phi^-(y)+\phi^-(x)\phi^-(y)$$
Supposing $x^0>y^0$, the difference between the time ordered and normal ordered
product is
$$T\left\{\phi(x)\phi(y)\right\}-N\left\{\phi(x)\phi(y)\right\} = [...]=\left[
\phi^-(x),\phi^+(y)\right]$$
and vice versa if $y^0>x^0$: the notation for this is the \textbf{contraction}
of the two fields
$$T\left\{\phi(x)\phi(y)\right\}-N\left\{\phi(x)\phi(y)\right\}:=
\overline{\phi(x)\phi(y)}:=
\begin{dcases*}
    \left[\phi^-(x),\phi^+(y)\right]\,\text{if } x^0>y^0\\
    \left[\phi^-(y),\phi^+(x)\right]\,\text{if } y^0>x^0
\end{dcases*}
$$
